% Source: Wald GR, physical PDF pp. 308--310
%         (printed pp. 295--297).
% Coverage: Chapter 11 problems 1--6.
% Figures/tables/footnotes: none.
% Status: source-reviewed and compile-clean.
% Uncertainties: none.

\chapterproblems

\begin{enumerate}[label=\arabic*.,leftmargin=2em,itemsep=1.2ex]
\item
Write out the coordinate components of the gauge condition
\(\tilde\nabla_a\tilde\nabla_b\Omega=0\) on \(\mathcal I^+\) in the coordinates
\((\Omega,u,\theta,\phi)\) introduced in the text. Show that these relations imply
that \(\tilde g_{uu}\), \(\tilde g_{u\theta}\), and \(\tilde g_{u\phi}\) are all
\(\order(\Omega^2)\) as \(\Omega\to0\).

\item
Let \((\Sigma,h_{ab},K_{ab})\) be an asymptotically flat initial data set. Introduce
coordinates \(X,Y,Z\) in a neighborhood of \(\Lambda\) such that the gradients of
these functions are orthonormal at \(\Lambda\). Show that the condition that
\(\tilde h_{ab}\) is \(C^{>0}\) at \(\Lambda\) implies that in these coordinates, the
unphysical metric components are \(\diag(1,1,1)+\order(R)\) as \(R\to0\), where
\(R=(X^2+Y^2+Z^2)^{1/2}\), and that the first derivatives of the metric components
are \(\order(1)\) as \(R\to0\). Show further that condition (iii) of the definition of
an asymptotically flat initial data set implies
\(\Omega=R^2[1+\order(R)]\). Introduce the coordinates
\(x=X/R^2\), \(y=Y/R^2\), \(z=Z/R^2\) in the physical spacetime and show that the
physical metric components take the form \(\diag(1,1,1)+\order(1/r)\) as
\(r\to\infty\) where \(r=(x^2+y^2+z^2)^{1/2}\). Show further that the first
coordinate derivatives of these components are \(\order(1/r^2)\) as \(r\to\infty\).
Finally, show that condition (iv) implies that the components of \(K_{ab}\) and the
physical Ricci tensor \({}^{(3)}R_{ab}\) in this asymptotically Euclidean coordinate
system are, respectively, \(\order(1/r^2)\) and \(\order(1/r^3)\).

\item
\begin{enumerate}[label=\alph*),leftmargin=2em,itemsep=.45ex]
\item Show that the Schwarzschild solution is asymptotically flat at future null
infinity.
\item Show that every \(t=\text{constant}\) hypersurface in the Schwarzschild
solution is an asymptotically flat initial data set.
\end{enumerate}
\noindent
(In fact, the Schwarzschild spacetime satisfies all the conditions of our definition
of asymptotic flatness at spatial and null infinity, not just the above separate
properties. For a proof, see Ashtekar and Hansen 1978.)

\item
Consider a stationary solution with stress-energy \(T_{ab}\) in the context of
linearized gravity (see Section~4.4). Choose a global inertial coordinate system for
the flat metric \(\eta_{ab}\) so that the \lq\lq time direction\rq\rq\
\((\partial/\partial t)^a\) of this coordinate system agrees with the stationary
Killing field to zeroth order.
\begin{enumerate}[label=\alph*),leftmargin=2em,itemsep=.45ex]
\item Show that the conservation equation, \(\partial^aT_{ab}=0\), implies
\(\int T_{\mu\nu}\,d^3x=0\), where the integral is taken over a \(t=\text{constant}\)
slice and \(\mu\) and \(\nu\) take any values except \(\mu=\nu=0\) (where
\(x^0=t\)).
\item Show, therefore, that to first order in deviation from flatness, the general
formula \eqref{eq:11.2.10} for \(M\) in a stationary spacetime reduces to
\(M=\int T_{00}\,d^3x\), i.e., we obtain the usual formula of special relativity.
Note, however, that the Komar integral \eqref{eq:11.2.9} taken over a finite sphere
lying within the matter distribution does not, in general, equal the mass contained
within that sphere.
\end{enumerate}

\item
\begin{enumerate}[label=\alph*),leftmargin=2em,itemsep=.45ex]
\item Calculate the right side of Eq.~\eqref{eq:11.2.9} for the Schwarzschild
solution and show that we get the Schwarzschild mass parameter \(M\). Thus, for a
static, spherically symmetric fluid star we have two formulas for \(M\), namely
Eq.~\eqref{eq:6.2.10} and Eq.~\eqref{eq:11.2.10}.
\item Show that the equality of the two formulas \eqref{eq:6.2.10} and
\eqref{eq:11.2.10} for \(M\) yields the relation
\[
  \int_\Sigma(\rho+3p)e^\phi\,dV
  =\int_\Sigma\rho\left(1-\frac{2m(r)}{r}\right)^{1/2}\,dV,
\]
where the integral is taken over a hypersurface \(\Sigma\) orthogonal to the static
Killing field \(\xi^a\), \(dV\) denotes the natural volume element on \(\Sigma\), and
\(\phi\equiv\tfrac12\ln(-\xi_a\xi^a)\) is the general relativistic analog of the
Newtonian gravitational potential (see Section~6.2).
\item Put in the \(G\)'s and \(c\)'s in the formula of part (b) and examine the
Newtonian limit, \(c\to\infty\). In this limit both sides of the above equation reduce
to the rest mass of the fluid. Show that to the next nonvanishing approximation in
\(1/c\), the equation of part (b) yields
\[
  E_g=-2K,
\]
where \(E_g\) is given by the usual Newtonian formula for gravitational potential
energy, and \(K\equiv\tfrac32\int_\Sigma p\,dV\), so that in the case of a monatomic
ideal gas, \(K\) is just the total thermal energy. Thus, the equation of part (b) may
be viewed as the general relativistic version of the Newtonian virial theorem.
\end{enumerate}

\item
In an axisymmetric, asymptotically flat spacetime with axial Killing field \(\psi^a\),
we may define the total angular momentum, \(J\), by
\[
  J=\frac{1}{16\pi}\int_S\epsilon_{abcd}\nabla^c\psi^d,
\]
where the integral is taken over a sphere, \(S\), in the asymptotic region where
\(T_{ab}\) is assumed to vanish. (Note the factor of \(-2\) difference between this
formula and the formula \eqref{eq:11.2.9} for \(M\).)
\begin{enumerate}[label=\alph*),leftmargin=2em,itemsep=.45ex]
\item Parallel the arguments given in Section~11.2 for \(M\) to show that \(J\) is
independent of \(S\) and that \(J\) is given by
\[
  J=-\int_\Sigma T_{ab}n^a\psi^b\,dV,
\]
where the hypersurface \(\Sigma\) is chosen so that \(\psi^a\) is tangent to
\(\Sigma\). Note that the fact that \(J\) is independent of \(S\) can be interpreted as
saying that in an axisymmetric spacetime angular momentum cannot be radiated away
by gravitational radiation.
\item Show that \(J=0\) in any static, axisymmetric spacetime, i.e., a spacetime
which possesses a hypersurface orthogonal timelike Killing field \(\xi^a\) with
\(\xi^a\psi_a=0\).
\item Calculate \(J\) for the charged Kerr metric (Eq.~\eqref{eq:12.3.1} of
Chapter~12) and show that \(J=Ma\).
\end{enumerate}
\noindent
(For discussions of the definition of angular momentum in nonaxisymmetric
asymptotically flat spacetimes, see Ashtekar 1980, Ashtekar and Streubel 1981,
Geroch and Winicour 1981, and Ashtekar and Winicour 1982.)
\end{enumerate}
